\documentclass[11pt]{article}
\usepackage[margin=1in]{geometry}
\usepackage[T1]{fontenc}
\usepackage[utf8]{inputenc}
\usepackage{hyperref}

\title{AgentFlow -- No-Code Agent Framework Based on Logical Primitives}
\author{
  Dmitry Lande\textsuperscript{1}(ORCID 0000-0003-3945-1178)
  \and
  Leonard Strashnoy\textsuperscript{2}(ORCID 0009-0008-5575-0286)
}
\date{June 2025}

\begin{document}
\maketitle

\begin{center}
{\small
\textsuperscript{1} National Technical University of Ukraine -- Igor Sikorsky Kyiv Polytechnic Institute\\
\textsuperscript{2} University of California, Los Angeles (UCLA)\\[0.5em]
SSRN: \url{https://papers.ssrn.com/sol3/papers.cfm?abstract_id=5285664}
}
\end{center}

\begin{abstract}
This article presents the AgentFlow framework -- a new paradigm of no-code
programming built on natural-language structural pseudocode, executed through Large Language
Models (LLMs). The framework utilizes logical primitives -- Condition, Loop, Function, Label,
Goto -- as fundamental building blocks for constructing complex agent behaviors. An agent
model is introduced, where each agent has: a role, state, logic, and communication capability.
The concept of a swarm of agents is proposed, enabling LLMs to simulate parallel task execution
and process complex queries without writing code. The article also describes a memory system
for self-learning, provides examples of data analysis, and outlines formal rules for agent-to-agent
interaction.
\end{abstract}

\noindent\textbf{Keywords:} no-code programming, large language models (LLM), agent model, swarm of agents,
logical primitives, parallel execution via prompts, agent self-learning

\section{Introduction}
% Port body from published/ssrn-5285664.pdf as you revise.

\section{Methodology}
\subsection{Basic primitives of no-code programming}
\subsection{Agent model}
\subsection{Swarm of agents}
\subsection{Memory and self-learning}

\section{Results}

\section{Discussion}

\section{Conclusions}

\bibliographystyle{plain}
\bibliography{references}

\end{document}
